\label{sec:introduction}

The citizens of Waukesha County, Wisconsin testified before the Wisconsin
Redistricting Commission in 2021 that their communities belonged in the same
legislative district as the other Milwaukee suburban counties to the north
and east.
Their argument was not about geography alone---Waukesha is not adjacent to
all the communities they claimed to share a district with.
Their argument was about daily life: tens of thousands of Waukesha residents
drive the same freeways, park in the same lots, and work in the same office
parks in Milwaukee's western suburbs.
They were, they testified, part of the same economic community.
The Commission heard the testimony, found it compelling, but had no quantitative
tool to measure the strength of the commuting community bond.

This paper provides that measurement.

The commuting shed concept has a long history in regional science.
\citet{tolbert1996labor} defined functional economic areas on the basis of
inter-county commuting flows, arguing that the daily commuting sphere---not
the municipal or county boundary---is the appropriate unit of economic community.
This work has been applied to labor market delineation, retail trade area
analysis, and economic development planning.
But it has not been applied to redistricting, where the relevant question is
not ``what is the optimal functional economic area size?'' but ``do these two
adjacent tracts belong to the same economic community?''

The M-track community character framework, introduced in M.0
\citep{m0framework}, operationalizes diverse community signals as edge
weights in the census-tract adjacency graph.
M.4 is the commuting shed implementation paper.
It uses the Census Bureau's LODES Origin-Destination (OD) file---which records
home-block to work-block commute flows for all UI-covered employees---to
compute a per-tract commuting destination distribution and measure the
similarity between adjacent tracts' distributions using the weighted
Jaccard coefficient.

Commuting shed similarity is the right M-track signal for polycentric
metropolitan areas for several reasons:

\begin{enumerate}
  \item \textbf{Functional community, not just residential proximity.}
    Two adjacent census tracts may be geographically contiguous but serve
    completely different employment centers.
    Tract A's workers commute north to a technology corridor; tract B's workers
    commute south to a port logistics hub.
    Geographic proximity implies no community bond; commuting destination
    overlap captures the functional difference.
  \item \textbf{Revealed preference.}
    Unlike survey-based community identification or administrative zone
    membership, commute flows are behavioral data.
    They reflect what workers actually do, not what planners designate or
    what residents aspire to.
    The commuting shed is the economic community as expressed through daily
    behavior.
  \item \textbf{Complementary to M.1.}
    M.1 (economic character, \citep{m1econchar}) weights tracts by their
    employment profile: where jobs are located.
    M.4 weights tracts by their residential commuting profile: where workers
    from each tract actually go.
    M.1 identifies employment centers; M.4 identifies their residential
    catchment areas.
    Together they provide a complete picture of the economic community.
  \item \textbf{Partisan neutrality.}
    LODES OD data contains no electoral, racial, or demographic information.
    The commute flow records are administrative data derived from state
    UI wage records matched to employer geocodes.
\end{enumerate}

\paragraph{Distinction from M.1 economic character.}
A high-employment tract (M.1 jobs per resident $> 5$) represents an employment
center.
M.4 identifies which residential tracts send workers to that center.
Two residential tracts that both send workers primarily to the same tech park
have a commuting-shed community bond; two residential tracts that send workers
to different employment centers do not---even if they are adjacent.
The M.4 signal fires where M.1 is silent: in residential areas with low JPR
values that cannot be distinguished by economic character but can be
distinguished by the direction of their daily commute.

\paragraph{Paper organization.}
Section~\ref{sec:background} describes the LODES OD data structure and
the functional economic area concept.
Section~\ref{sec:algorithm} defines the commuting destination distribution,
weighted Jaccard similarity, sparsity handling, and edge weight formula.
Section~\ref{sec:results} describes the experimental design and defers
quantitative results to Phase~2.
Section~\ref{sec:legal} provides a legal usage guide including airport
communities, port communities, and economic commuting shed recognition.
Section~\ref{sec:conclusion} concludes.
